\documentclass[%
  % aspectratio=169,
  aspectratio=1610,
]{beamer}

%-----------------------------------------
% Set Theme Options
%-----------------------------------------
\usetheme[%
  logo=magicc, % byu|magicc|radlab
  footline=default, % default|minimal|none
  nav=top, % top|bottom|none
  framenumber=current, % current|total
]{byu}

% (Optional) Scale height of frametitle bar: default is 100
\renewcommand{\byuFrametitleBarScalePercent}{125}

%%%% Set notes style
\setbeameroption{show notes on second screen=right}
\setbeamertemplate{note page}[byu] % built-ins are default|compress|plain
%% (fine tuned control of byu notes theme)
% \byuNotesPreviewOff
% \byuNotesContextOff
% \byuNotesTitleOff
% \byuNotesRuleOff
% \byuNotesSlideNumberOff
% \byuNotesBodyBesidePreviewOn % controls whether notes are left or below of preview
% \setlength{\byuNotesPageMargin}{1ex}
% \setlength{\byuNotesColumnGap}{1ex}
% \setlength{\byuNotesPreviewPadding}{1ex}
% \renewcommand{\byuNotesPreviewScale}{0.3}

%-----------------------------------------
% Import packages
%-----------------------------------------
% Font packages for XeLaTeX or LuaLaTeX (pdflatex will not work with these)
\usepackage{fontspec}
\usepackage[mathbf=sym]{unicode-math} % replaces deprecated \mathbf with new \symbf

%-----------------------------------------
% Choose fonts/typefaces
%-----------------------------------------
\setmainfont{Source Serif Pro}
\setsansfont{Source Sans Pro} % beamer mainly uses sans
\setmonofont{SauceCodePro Nerd Font Mono}

%%%% Math font options
% \setmathfont{Tex Gyre Pagella Math}
% \setmathfont{Latin Modern Math}
% \setmathfont{NewCMSansMath-Regular.otf}
\setmathfont{Noto Sans Math} % doesn't define U+2022 (bullet) so you need the next line too
\setmathfont{Noto Sans}[range="2022, Scale=MatchUppercase]

%--------------------------------------------
% Override bullet points (optional)
%--------------------------------------------

%------ Fun bullet points (may require a nerd font)
% \newfontfamily\emojifont{SauceCodePro Nerd Font}
% \setbeamertemplate{itemize item}{\emojifont \scriptsize }
% \setbeamertemplate{itemize subitem}{\emojifont \scriptsize 🤖}
% \setbeamertemplate{itemize subsubitem}{\emojifont \scriptsize }

%------ Standard bullet points
% \setbeamertemplate{itemize subsubitem}[triangle]
% \setbeamertemplate{itemize subsubitem}[square]
% \setbeamertemplate{itemize subsubitem}[star]
% \setbeamertemplate{itemize subsubitem}[ball]
% \setbeamertemplate{itemize item}{$\circ$}
% \setbeamertemplate{itemize subitem}{$\diamond$}

%------------------------------------------------
% This adds Table of Contents before each section
%------------------------------------------------
\setcounter{tocdepth}{3}
\AtBeginSection[]{% If you want to show the outline before each section
  \begin{frame}{Outline}
    \tableofcontents[
      currentsection,
      sectionstyle=show/shaded,              % Gray out non-active sections
      subsectionstyle=show/show/hide,        % Hide all subsections
      subsubsectionstyle=hide/hide/hide/hide % Hide all subsubsections
    ]
  \end{frame}
}
\AtBeginSubsection[]{% If you want to show the outline before each subsection
  \begin{frame}{Outline}
    \tableofcontents[
      currentsection,
      currentsubsection,
      % sectionstyle=shaded/shaded,            % Gray out all sections
      sectionstyle=show/shaded,              % Gray out non-active sections
      subsectionstyle=show/shaded/hide,      % Highlight active subsection, hide others
      subsubsectionstyle=show/show/hide/hide % Reveal subsubsections ONLY for active branch
    ]
  \end{frame}
}

%--------------------------------------------
% Custom commands
%--------------------------------------------
\newcommand{\speak}[2]{% for use inside of note environments
  \textbf{\textcolor{byuRoyal}{#1:}} #2 \par \vspace{0.5em}
}
\newcommand{\takeaway}[1]{

}

%--------------------------------------------
% Front matter
%--------------------------------------------
\title[Short Title]{Super Awesome Long Title}
\institute[BYU]{Brigham Young University}
\author[Short Author]{Long Author John Doe}
\date[Jun 2026]{June 24, 2026}


%--------------------------------------------
% Main content
%--------------------------------------------
\begin{document}

%--------------------------------------------------
\begin{frame}[plain,noframenumbering]
\titlepage
\end{frame}

%--------------------------------------------------
\begin{frame}[plain,noframenumbering]
\centering
\resizebox{\textheight}{!}{%
\begin{tikzpicture}
  \pic[byu color=byuNavy] {byu_logo};
\end{tikzpicture}%
}
\end{frame}

%--------------------------------------------------
\begin{frame}[plain,noframenumbering]
\centering
\resizebox{\textheight}{!}{%
\begin{tikzpicture}
  \pic[
    magicc chevron color=byuRoyal,
    magicc arc color=magiccGray,
    magicc text color=magiccDarkGray,
    byu color=byuNavy,
  ] {magicc_logo};
\end{tikzpicture}%
}
\end{frame}

%--------------------------------------------------
\begin{frame}[plain,noframenumbering]
\centering
\resizebox{0.72\textheight}{!}{%
\begin{tikzpicture}
  \pic[
    outline=radlabBlue,
    background=none,
    text box=radlabBlue,
    text=white,
    text claws=white,
    base=radlabBlue,
    torso=radlabBlue,
    upper arms=black,
    forearms=radlabBlue,
    hands=black,
    face=radlabBlue,
    eye=white,
    helmet=black,
    joints=white,
  ] {radlab_logo};
\end{tikzpicture}%
}
\end{frame}


%--------------------------------------------------
% Table of Contents
%--------------------------------------------------
\begin{frame}{Outline}
  \tableofcontents[subsectionstyle=hide,subsubsectionstyle=hide]
\end{frame}


%--------------------------------------------------
\section{Introduction}
%--------------------------------------------------


%--------------------------------------------------
\begin{frame}{Overview}

Here is a test slide:

\begin{itemize}
\item Background
\item Methods
\item Results
\item Takeaways
\end{itemize}

\note[item]{This is my note}

\end{frame}

\subsection{List Tests}
\subsubsection{Combined}
%--------------------------------------------------
\begin{frame}{List Tests}

\begin{columns}
  \begin{column}{0.5\textwidth}
    \begin{itemize}
      \item Level 1
      \begin{itemize}
        \item Level 2
        \begin{itemize}
          \item Level 3
          \item This is the max level depth
        \end{itemize}
      \end{itemize}
    \end{itemize}
  \end{column}

  \begin{column}{0.5\textwidth}
    \begin{enumerate}
      \item Level 1
      \begin{enumerate}
        \item Level 2
        \begin{enumerate}
          \item Level 3
          \item This is the max level depth
        \end{enumerate}
      \end{enumerate}
    \end{enumerate}
  \end{column}
\end{columns}

\end{frame}

\subsubsection{Bullet}
%--------------------------------------------------
\begin{frame}{Bulleted Lists}

\begin{itemize}
  \item Level 1
  \begin{itemize}
    \item Level 2
    \begin{itemize}
      \item Level 3
      \item This is the max level depth
    \end{itemize}
  \end{itemize}
\end{itemize}

\end{frame}

\subsubsection{Numbered}
%--------------------------------------------------
\begin{frame}{Numbered Lists}

\begin{enumerate}
  \item Level 1
  \begin{enumerate}
    \item Level 2
    \begin{enumerate}
      \item Level 3
      \item This is the max level depth
    \end{enumerate}
  \end{enumerate}
\end{enumerate}

\end{frame}

\subsection{Block Tests}
%--------------------------------------------------
\begin{frame}{Blocks}


\begin{block}{Block}
  \begin{quote}
    It's not the plane, it's the pilot. --Rooster
  \end{quote}
\end{block}

\begin{alertblock}{Alert Block}
  \begin{quote}
    Trust your instincts. Don't think. Just do. --Maverick
  \end{quote}
\end{alertblock}

\begin{example}
  \begin{quote}
    You can be my wingman anytime. --Iceman
  \end{quote}
\end{example}

\end{frame}


%--------------------------------------------------
\section{Background}
%--------------------------------------------------


%--------------------------------------------------
\begin{frame}{Related Work}
  \centering
  {\LARGE Sources}
\end{frame}


%--------------------------------------------------
\section{Methods}
%--------------------------------------------------


%--------------------------------------------------
\begin{frame}{Math}
  \begin{equation}
    \int_{-\infty}^{\infty} e^{-x^2} dx = \sqrt{\pi}
  \end{equation}

  \begin{equation}
    \min_{\mathbf{u}} \quad \sum_{k=0}^{T} \| \tilde{\mathbf{x}}_k \|_Q^2 + \| \tilde{\mathbf{u}}_k \|_R^2
  \end{equation}

  % \begin{equation}
  %   \mathbf{u}_{\mathrm{min}} \leq \mathbf{u}_k \leq \mathbf{u}_{\mathrm{max}}
  % \end{equation}
  \begin{align}
    \mathbf{u}_{\mathrm{min}} \leq \mathbf{u}_k \leq \mathbf{u}_{\mathrm{max}} \\
    \mathbf{u}_{\textrm{min}} \leq \mathbf{u}_k \leq \mathbf{u}_{\textrm{max}} \\
    \mathbf{u}_{\textnormal{min}} \leq \mathbf{u}_k \leq \mathbf{u}_{\textnormal{max}}
  \end{align}

  \note{
    \speak{Hook}{This is the square root of pi!}
    \speak{Key Point}{Regularizing cost function}
    \speak{Transition}{Moving on to a code example}
  }
\end{frame}


%--------------------------------------------------
\begin{frame}[fragile]{Code}
\begin{scriptsize} % mono font is quite large, so we shrink it down a bit
\begin{verbatim}
def fibonacci(n: int) -> list[int]:
    values = [0, 1]
    while len(values) < n:
        values.append(values[-1] + values[-2])
    return values

print(fibonacci(10))  # -> [0, 1, 1, 2, 3, 5, 8, 13, 21, 34]
\end{verbatim}
\end{scriptsize}
\end{frame}

%--------------------------------------------------
\begin{frame}{Methods}
  \centering
  {\LARGE Do this and you're cool}
\end{frame}

%--------------------------------------------------
\begin{frame}{More Content}
  \centering
  {\LARGE More cool stuff}
\end{frame}


%--------------------------------------------------
\section{Results}
%--------------------------------------------------


%--------------------------------------------------
\begin{frame}{Results}
  \centering
  {\LARGE Awesome stuff}
\end{frame}

%--------------------------------------------------
\begin{frame}{More Content}
  \centering
  {\LARGE More cool stuff}
\end{frame}


%--------------------------------------------------
\section{Takeaways}
%--------------------------------------------------


%--------------------------------------------------

\begin{frame}{Conclusion}
  \centering
  {\LARGE This is why I should graduate}
\end{frame}


\end{document}
